\section{Recomendaciones}
\label{sec:recomendaciones}

\begin{table}[htbp]
    \centering
    \caption{Recomendaciones derivadas del an\'{a}lisis de la l\'{i}nea SIM-002.}
    \label{tab:recomendaciones}
    \setcounter{itemcount}{0}
    \begin{tabularx}{\textwidth}{@{}NlX@{}}
        \toprule
        \multicolumn{1}{c}{\textbf{Item}} & \textbf{Categoria} & \textbf{Recomendacion} \\
        \midrule
        & Construcci\'{o}n & Construir la soporter\'{i}a estrictamente conforme al isom\'{e}trico PCF103: los soportes r\'{i}gidos verticales (+Y) de los nodos 25 y 70, la restricci\'{o}n r\'{i}gida del nodo 130 y las holguras de los soportes de los extremos (nodos 10 y 170) son parte integral del comportamiento verificado; cualquier cambio de posici\'{o}n, tipo u holgura de soporte invalida el presente an\'{a}lisis. \\
        & Ingenier\'{i}a & Remitir las cargas de la restricci\'{o}n del nodo 10 (Tablas~\ref{tab:restricciones-ope} a~\ref{tab:restricciones-exp}) al responsable de la bomba de recirculaci\'{o}n 1 DDW para confirmar la aceptaci\'{o}n de las cargas en la boquilla de descarga. \\
        & Documentaci\'{o}n & Conservar el job CAESAR P2603-PR-SIM-002 y el PCF corregido como registro del an\'{a}lisis; ante cualquier modificaci\'{o}n de la l\'{i}nea (di\'{a}metros, recorrido, condiciones de operaci\'{o}n), re-ejecutar el modelo y re-emitir este informe con una nueva revisi\'{o}n. \\
        \bottomrule
    \end{tabularx}
\end{table}
